\section{Event Sources}

TODO

\subsection{Static Events}

\subsubsection*{Tracepoints}



Aka. static event tracing.

There are more than 1000 traceoints available (\verb|grep -R TRACE_EVENT include/trace/events|)

USEFUL: \url{https://mmi.hatenablog.com/entry/2018/03/04/052249}

Interesting: static-key mechanism
\url{https://www.kernel.org/doc/html/latest/staging/static-keys.html}

Put an entire implementation of a simple module that registers a
tracepoint at the appendix.

\subsubsection*{Raw Tracepoints}

TODO

\subsubsection*{Syscall Tracepoints}

TODO

\subsubsection*{User Statically-Defined Tracing}

(USDT)
Note: this is actually a uprobe in disguise, it expects to attach to
tracepoints defined in userspace programs.

TODO: write some code to understand this

\subsection{Dynamic Events}

TODO

\subsubsection*{Kprobes and Kretprobes}

TODO


Architecture
------------

There are kprobes and kretprobes. From the docs:

```
When a kprobe is registered, Kprobes makes a copy of the probed
instruction and replaces the first byte(s) of the probed instruction
with a breakpoint instruction (e.g., int3 on i386 and x86\_64).

When a CPU hits the breakpoint instruction, a trap occurs, the CPU’s
registers are saved, and control passes to Kprobes via the
\verb|notifier_call_chain| mechanism. Kprobes executes the \verb|pre_handler|
associated with the kprobe, passing the handler the addresses of the
kprobe struct and the saved registers.

Next, Kprobes single-steps its copy of the probed instruction. (It
would be simpler to single-step the actual instruction in place, but
then Kprobes would have to temporarily remove the breakpoint
instruction. This would open a small time window when another CPU
could sail right past the probepoint.)

After the instruction is single-stepped, Kprobes executes the
\verb|post_handler|, if any, that is associated with the kprobe. Execution
then continues with the instruction following the probepoint.
```

Picture: show this graphically. Take inspiration to the lwn article.

Kretprobes do the same but by playing with the return address, which
is places pushed onto the stack by the caller.

Note that kprobes can do everything, which is very dangerous.

Kprobes may perform jump optimization if the architecture supports
it and the kernel was built with \verb|CONFIG_OPTPROBES=y|. In this case,
instead of a breakpoint instruction, the CPU performs a jump instead.

Note that kprobes do not work well with inlining.

Interface
---------

API is in include/linux/kprobes.h

\begin{verbatim}

    struct kprobe { ... }
    typedef int (*kprobe_pre_handler_t) (struct kprobe *, struct pt_regs *);
    typedef void (*kprobe_post_handler_t) (struct kprobe *, struct pt_regs *,
		    		       unsigned long flags);

\end{verbatim}

Kprobes have probe handlers, a kprobe manager and architecture
dependent layers for messing with instruction and exceptions.

You call register / undegister kprobe / kretprobe. You can enable and
disable them.

You can use a symbol instead of an address, the kernel will resolve
the address from the symbol using System.map symbol table.

The list of registered kprobes is visible under DEBUGFS/kprobes/list


Implementation
--------------

kernel/kprobes.c implements the mangaer, which stores an hashmap of
kprobes indexed by their address at which the probe is placed, and
blacklisted entries. Insertion and deletion to this table is guarded
by a mutex to prevent these operations from executing simultaneously.

\begin{verbatim}

Whenever a probe is hit, the probe handler is called with interrupts disabled. Interrupts are disabled because handling a probe is a multiple step process which involves breakpoint handling and single-step execution of the probed instruction. There is no easy way to save the state between these operations hence interrupts are kept disabled during probe handling.

\end{verbatim}

Example implementation usage \verb|kprobe_table|:

kernel/kprobes.c:378

\begin{verbatim}

struct kprobe *get_kprobe(void *addr)
{
	struct hlist_head *head;
	struct kprobe *p;

	head = &kprobe_table[hash_ptr(addr, KPROBE_HASH_BITS)];
	hlist_for_each_entry_rcu(p, head, hlist,
				 lockdep_is_held(&kprobe_mutex)) {
		if (p->addr == addr)
			return p;
	}

	return NULL;
}

\end{verbatim}

Taking x86\_64 as an example, the interrupt handler for int3 is in
arch/x86/kernel/traps.c and indeed calls the kprobe handler if kprobes
are enabled:

\begin{verbatim}

static bool do_int3(struct pt_regs *regs)
{
	int res;
  // ...
#ifdef CONFIG_KPROBES
	if (kprobe_int3_handler(regs))
		return true;
#endif
  // ...
}

\end{verbatim}

Which is defined in arch/x86/kernel/kprobes/core.c. Show the call
stack here

Main source of information:
\url{https://www.kernel.org/doc/html/latest/trace/kprobes.html}
COOL: \url{https://lwn.net/Articles/132196/}

\subsubsection*{Function Hooks}

TODO

Technology: gcc -pg


\subsubsection*{Fprobe}

TODO

\subsubsection*{Uprobes and Uretprobes}

TODO

\subsubsection*{uprobe-multi}

TODO

\subsubsection*{Eprobes}

TODO

\subsection{Perf events}

TODO

\subsubsection*{Software Events}

TODO

\subsubsection*{Hardware Events}

TODO

Also talk about hardware tracing here. Perf events are easy, they just
notify that an even has appened. There are hardware events and
software events.

There is this abstraction of a performance mesuring unit. This is an
generic object that can manages a set of performance events and
counters. It handles registration, stopping and reading of events.
Code is in \verb|linux/include/perf_event.h|; struc pmu: 329.


Security implications:
\url{https://www.kernel.org/doc/html/latest/admin-guide/perf-security.html}

\begin{verbatim}

perf_event_open(2)

       A call to perf_event_open() creates a file descriptor that  allows  mea‐
       suring performance information.  Each file descriptor corresponds to one
       event  that is measured; these can be grouped together to measure multi‐
       ple events simultaneously.

       Events can be enabled and disabled in two ways:  via  ioctl(2)  and  via
       prctl(2).  When an event is disabled it does not count or generate over‐
       flows but does continue to exist and maintain its count value.

       Events  come  in two flavors: counting and sampled.  A counting event is
       one that is used for counting the aggregate number of events that occur.
       In general, counting event results are gathered with a read(2) call.   A
       sampling  event  periodically  writes  measurements to a buffer that can
       then be accessed via mmap(2).

\end{verbatim}

Another resource:
  \url{https://github.com/torvalds/linux/blob/master/tools/perf/design.txt}
